\documentclass[11pt]{article}

\usepackage[a4paper,margin=1in]{geometry}
\usepackage{amsmath,amssymb,amsthm,mathtools}
\usepackage[T1]{fontenc}
\usepackage{lmodern}
\usepackage{microtype}
\usepackage{enumitem}
\usepackage[colorlinks=true,linkcolor=blue,citecolor=blue,urlcolor=blue]{hyperref}

\newtheorem{theorem}{Theorem}
\newtheorem{proposition}{Proposition}
\newtheorem{lemma}{Lemma}
\newtheorem{corollary}{Corollary}
\theoremstyle{remark}
\newtheorem{remark}{Remark}

\newcommand{\Z}{\mathbb{Z}}

\title{A Self-Contained Analysis of the Integer Equation\\[0.3em]
\large \texorpdfstring{$y(x^3-z^2)=z+1$}{y(x^3-z^2)=z+1}}
\author{}
\date{}

\begin{document}
\maketitle

\begin{abstract}
We give a completely elementary proof that the Diophantine equation
\[
        y(x^3-z^2)=z+1
\]
has infinitely many integer solutions.  More precisely, the entire fibre with
$z=-1$ is determined explicitly: it is the union of the two infinite families
\[
        (x,y,z)=(t,0,-1) \quad (t\in\Z)
\]
and
\[
        (x,y,z)=(1,t,-1) \quad (t\in\Z).
\]
For completeness, we also record a full structural characterization of all
integer solutions by a single divisibility condition, together with the unique
exceptional zero-denominator case.
\end{abstract}

\section{The problem}

We study integer triples $(x,y,z)\in\Z^3$ satisfying
\begin{equation}\label{eq:main}
        y(x^3-z^2)=z+1.
\end{equation}
The requested alternatives are either to list all integer solutions or to prove
that infinitely many integer solutions exist.  The latter is immediate once the
fibre $z=-1$ is examined carefully.  The point of this note is to present the
argument in a fully explicit form, with all elementary facts used in the proof
stated and proved.

\section{Elementary preliminaries}

\begin{lemma}[Zero-product property in the integers]\label{lem:zpp}
Let $a,b\in\Z$.  If $ab=0$, then $a=0$ or $b=0$.
\end{lemma}

\begin{proof}
Suppose, for contradiction, that $a\ne 0$ and $b\ne 0$.  Then
$|a|\ge 1$ and $|b|\ge 1$, so
\[
        |ab|=|a|\,|b|\ge 1.
\]
Hence $ab\ne 0$, contradicting the hypothesis $ab=0$.  Therefore at least one
of $a,b$ is zero.
\end{proof}

\begin{lemma}[The integer cube equal to one]\label{lem:cubeone}
Let $x\in\Z$.  If $x^3=1$, then $x=1$.
\end{lemma}

\begin{proof}
By trichotomy, exactly one of the following holds: $x\le 0$, $x=1$, or
$x\ge 2$.  If $x\le 0$, then $x^3\le 0$, so $x^3\ne 1$.  If $x\ge 2$, then
$x^3\ge 2^3=8$, so again $x^3\ne 1$.  The only remaining possibility is
$x=1$.
\end{proof}

\section{The fibre \texorpdfstring{$z=-1$}{z=-1}}

\begin{proposition}[Exact determination of the fibre $z=-1$]\label{prop:fibre}
The integer solutions of \eqref{eq:main} with $z=-1$ are exactly
\[
        \{(t,0,-1):t\in\Z\}\ \cup\ \{(1,t,-1):t\in\Z\}.
\]
The union is not disjoint, since both families contain $(1,0,-1)$.
\end{proposition}

\begin{proof}
Set $z=-1$ in \eqref{eq:main}.  Since $(-1)^2=1$ and $(-1)+1=0$, the equation
becomes
\begin{equation}\label{eq:zminusone}
        y(x^3-1)=0.
\end{equation}
By Lemma~\ref{lem:zpp}, equation \eqref{eq:zminusone} holds if and only if
\[
        y=0 \qquad \text{or} \qquad x^3-1=0.
\]
The second condition is $x^3=1$, which by Lemma~\ref{lem:cubeone} is equivalent
to $x=1$.

Thus a triple with $z=-1$ solves \eqref{eq:main} if and only if either
$y=0$, with $x$ arbitrary, or $x=1$, with $y$ arbitrary.  This is precisely
\[
        \{(t,0,-1):t\in\Z\}\ \cup\ \{(1,t,-1):t\in\Z\}.
\]
Conversely, every triple in this union satisfies \eqref{eq:zminusone}, and
therefore satisfies \eqref{eq:main} with $z=-1$.
\end{proof}

\section{A full structural characterization}

The preceding proposition already proves infinitude.  The following exact
criterion is included to remove any possible ambiguity about the role of the
factor $x^3-z^2$.

\begin{proposition}[Divisibility description of all integer solutions]\label{prop:structure}
An integer triple $(x,y,z)\in\Z^3$ satisfies
\[
        y(x^3-z^2)=z+1
\]
if and only if it belongs to exactly one of the following two classes:
\begin{enumerate}[label=\textup{(\roman*)}]
    \item $x^3-z^2\ne 0$, the integer $x^3-z^2$ divides $z+1$, and
    \[
            y=\frac{z+1}{x^3-z^2};
    \]
    \item $(x,z)=(1,-1)$ and $y$ is arbitrary.
\end{enumerate}
\end{proposition}

\begin{proof}
Put
\[
        D=x^3-z^2.
\]
Then the equation is
\begin{equation}\label{eq:D}
        yD=z+1.
\end{equation}

First assume $D\ne 0$.  If $(x,y,z)$ is a solution, then \eqref{eq:D} says
that $z+1$ is an integer multiple of $D$, namely the multiple $y$.  Hence
$D\mid (z+1)$ and
\[
        y=\frac{z+1}{D}=\frac{z+1}{x^3-z^2}.
\]
Conversely, if $D\ne 0$, $D\mid(z+1)$, and
$y=(z+1)/D$, then multiplying both sides by $D$ gives $yD=z+1$, which is
exactly the original equation.

It remains to handle $D=0$.  In that case the left-hand side of
\eqref{eq:D} is $0$, so \eqref{eq:D} is equivalent to
\[
        z+1=0.
\]
Thus $z=-1$.  Since $D=0$, we also have
\[
        x^3-z^2=0.
\]
Substituting $z=-1$ gives
\[
        x^3-1=0,
\]
so $x^3=1$, and Lemma~\ref{lem:cubeone} gives $x=1$.  Therefore the only
solutions with $D=0$ have $(x,z)=(1,-1)$ and arbitrary $y$.
Conversely, if $(x,z)=(1,-1)$, then $x^3-z^2=1-1=0$ and $z+1=0$, so
\[
        y(x^3-z^2)=y\cdot 0=0=z+1
\]
for every $y\in\Z$.

The two classes are disjoint because class (i) requires $x^3-z^2\ne 0$, while
class (ii) has $x^3-z^2=0$.
\end{proof}

\section{Infinitely many integer solutions}

\begin{theorem}\label{thm:infinite}
The Diophantine equation
\[
        y(x^3-z^2)=z+1
\]
has infinitely many integer solutions.
\end{theorem}

\begin{proof}
By Proposition~\ref{prop:fibre}, every triple of the form
\[
        (x,y,z)=(1,t,-1), \qquad t\in\Z,
\]
is a solution.  Directly substituting also verifies this without reference to
any previous result:
\[
        t(1^3-(-1)^2)=t(1-1)=0=(-1)+1.
\]
If $t_1,t_2\in\Z$ and $t_1\ne t_2$, then
\[
        (1,t_1,-1)\ne (1,t_2,-1),
\]
so these triples are pairwise distinct.  Since $\Z$ is infinite, the displayed
family contains infinitely many distinct integer solutions.
\end{proof}

\begin{corollary}
The equation also has the infinite solution family
\[
        (x,y,z)=(t,0,-1),\qquad t\in\Z.
\]
\end{corollary}

\begin{proof}
For every $t\in\Z$,
\[
        0\cdot(t^3-(-1)^2)=0=(-1)+1,
\]
so $(t,0,-1)$ is a solution.  Distinct values of $t$ give distinct triples.
\end{proof}

\begin{remark}
No finiteness classification is needed to answer the stated problem, because
Theorem~\ref{thm:infinite} proves one of the requested alternatives.  Still,
Proposition~\ref{prop:structure} gives an exact necessary-and-sufficient
condition for every integer solution: outside the unique exceptional case
$(x,z)=(1,-1)$, the value of $y$ is forced by divisibility.
\end{remark}

\end{document}
